% Bagian: first-order-logic
% Bab: beyond
% Subbagian: many-sorted-logic

\documentclass[../../../include/open-logic-section]{subfiles}

\begin{document}

\olfileid[id]{fol}{byd}{msl}

\olsection{Logika Banyak-Sorta}

Dalam logika orde pertama, variabel dan kuantor merentang atas satu
!!{domain}. Namun, sering kali berguna untuk mempunyai beberapa !!{domain}s
(yang saling lepas): misalnya, kita mungkin ingin mempunyai !!a{domain}
bilangan, !!a{domain} objek geometri, !!a{domain} fungsi dari bilangan ke
bilangan, !!a{domain} grup abelian, dan seterusnya.

Logika banyak-sorta menyediakan kerangka semacam ini. Kita mulai dengan suatu
daftar ``sorta''---``sorta'' suatu objek menunjukkan ``!!{domain}'' yang
seharusnya ditempatinya. Kemudian kita mempunyai !!{variable}s dan kuantor
untuk setiap sorta, serta (biasanya) !!{identity} untuk setiap sorta. Fungsi
dan relasi juga ``diberi tipe'' berdasarkan sorta objek yang dapat menjadi
argumennya. Selain itu, aturan-aturan logika orde pertama yang biasa tetap
digunakan, dengan versi aturan kuantor yang diulangi untuk setiap sorta.

Misalnya, untuk mengkaji hubungan internasional kita dapat memilih suatu
bahasa dengan dua sorta objek, yaitu warga negara Prancis dan warga negara
Jerman. Kita dapat mempunyai relasi uner ``minum anggur'' untuk objek sorta
pertama; relasi uner lain, ``makan sosis bratwurst'', untuk objek sorta kedua;
dan relasi biner ``membentuk pasangan menikah multinasional'' yang menerima
dua argumen, dengan argumen pertama dari sorta pertama dan argumen kedua dari
sorta kedua. Jika kita menggunakan variabel $a$, $b$, $c$ untuk merentang atas
warga negara Prancis dan $x$, $y$, $z$ untuk merentang atas warga negara
Jerman, maka
\[
\lforall[a][\lforall x][(\Atom{\Obj{MarriedTo}}{a,x} \lif
(\Atom{\Obj{DrinksWine}}{a} \lor \lnot \Atom{\Obj{EatsWurst}}{x})]]
\]
menyatakan bahwa jika seorang Prancis menikah dengan seorang Jerman, maka
orang Prancis itu minum anggur atau orang Jerman itu tidak makan sosis
bratwurst.

Logika banyak-sorta dapat ditanamkan ke dalam logika orde pertama dengan cara
yang alami, yakni dengan menggabungkan semua objek dari !!{domain}s
banyak-sorta menjadi satu !!{domain} orde pertama, menggunakan !!{predicate}s
uner untuk mencatat sorta, dan merelatifkan kuantor. Misalnya, bahasa orde
pertama yang bersesuaian dengan contoh di atas akan mempunyai !!{predicate}s
uner ``$\Obj{German}$'' dan ``$\Obj{French}$'', selain relasi-relasi lain yang
telah diuraikan, dengan persyaratan sorta dihapus. Kuantor bersorta
$\lforall[x][!A]$, dengan $x$ merupakan !!a{variable} dari sorta Jerman,
diterjemahkan menjadi
\[
\lforall[x][(\Atom{\Obj{German}}{x} \lif !A)].
\]
Kita perlu menambahkan aksioma yang memastikan bahwa sorta-sorta itu terpisah
---misalnya, $\lforall[x][\lnot (\Atom{\Obj{German}}{x} \land
  \Atom{\Obj{French}}{x})]$---serta aksioma yang menjamin bahwa ``minum
anggur'' hanya berlaku pada objek yang memenuhi predikat
$\Atom{\Obj{French}}{x}$, dan seterusnya. Dengan konvensi dan aksioma ini,
tidak sulit untuk menunjukkan bahwa !!{sentence}s banyak-sorta diterjemahkan
menjadi !!{sentence}s orde pertama, dan !!{derivation}s banyak-sorta
diterjemahkan menjadi !!{derivation}s orde pertama. Selain itu,
!!{structure}s banyak-sorta ``diterjemahkan'' menjadi !!{structure}s orde
pertama yang bersesuaian dan sebaliknya, sehingga kita juga memperoleh teorema
kelengkapan untuk logika banyak-sorta.

\end{document}
